\chapter{Declaración de uso responsable de la inteligencia artificial}
\label{anx:ia}

En cumplimiento de los principios de transparencia e integridad académica, y en
coherencia con lo indicado en la sección~\ref{sec:entorno-desarrollo}, el autor
declara haber empleado herramientas de inteligencia artificial generativa como
apoyo durante la realización de este \ac{TFG}.  La presente declaración detalla
de forma explícita la naturaleza, el alcance y los límites de dicho uso.

\section*{Finalidad de la declaración}

Esta declaración tiene por objeto documentar, de manera transparente, el uso de
herramientas de inteligencia artificial en el desarrollo del trabajo,
distinguiendo con claridad las tareas en las que dichas herramientas prestaron
apoyo de aquellas que constituyen aportación, criterio y responsabilidad
exclusivos del autor.

\section*{Herramientas empleadas}

Se ha utilizado el asistente conversacional Claude (Anthropic) como herramienta
de apoyo a lo largo del proyecto.  No se ha empleado ninguna otra herramienta de
inteligencia artificial generativa en la elaboración del sistema ni de esta
memoria.

\section*{Ámbitos de uso}

El apoyo de la herramienta se circunscribió a los siguientes ámbitos:

\begin{itemize}
  \item \textbf{Generación de código}, de forma principal en el \emph{frontend}
    y en las pruebas automatizadas, a partir de especificaciones y requisitos
    definidos por el autor.

  \item \textbf{Consultas técnicas y depuración}, como apoyo puntual en la
    resolución de dudas y en el diagnóstico de errores de integración entre los
    distintos servicios del sistema.

  \item \textbf{Revisión de la redacción de la memoria}, en aspectos de estilo,
    claridad, concisión y corrección, siempre sobre contenido elaborado por el
    autor y bajo su decisión y validación final.
\end{itemize}

\section*{Alcance y límites}

El uso de la herramienta se limitó en todo momento a un rol de asistencia, y
nunca de sustitución.  Fueron responsabilidad exclusiva del autor la concepción
del proyecto y la definición de sus objetivos, el diseño de la arquitectura y
del modelo de datos, la selección tecnológica, las decisiones de ingeniería, el
análisis de los resultados y la elaboración de las conclusiones.  Ningún
fragmento de código o de texto se incorporó al trabajo sin haber sido
previamente comprendido, revisado, verificado y, cuando fue necesario, corregido
o reescrito por el autor.

\section*{Verificación y validación}

Todo el material producido con apoyo de la herramienta fue sometido a
verificación.  El código se validó mediante su ejecución, la batería de pruebas
automatizadas y el análisis estático de calidad recogido en el
Anexo~\ref{anx:sonarqube}.  El texto se sometió a una revisión crítica orientada
a garantizar su exactitud, su coherencia con el sistema efectivamente
desarrollado y la ausencia de afirmaciones no respaldadas por el trabajo
realizado.

\section*{Declaración de responsabilidad}

El autor asume la plena y exclusiva responsabilidad sobre la totalidad del
contenido de esta memoria y del sistema desarrollado.  Las herramientas de
inteligencia artificial se han empleado como un apoyo a la productividad y al
aprendizaje, sin sustituir el criterio, el esfuerzo ni la comprensión propios, y
respetando en todo momento los principios de honestidad e integridad académica
que rigen la elaboración de un Trabajo Fin de Grado.


% Local Variables:
%  coding: utf-8
%  mode: latex
%  mode: flyspell
%  ispell-local-dictionary: "castellano8"
% End:
